% ======================================================================
%  Wenhan Cao - Curriculum Vitae
%  Compile with pdflatex (twice, for hyperlinks).
% ----------------------------------------------------------------------
%  Structure: each section uses small helper macros defined below, so
%  adding a publication = copying one \pubitem{...}{...}{...}{...} block.
% ======================================================================
\documentclass[11pt,a4paper]{article}

% ---------- packages ----------
\usepackage[a4paper,margin=0.63in,top=0.7in,bottom=0.7in]{geometry}
\usepackage[T1]{fontenc}
\usepackage[utf8]{inputenc}
\usepackage{charter}            % Bitstream Charter serif, close to Georgia; swap for \usepackage{mathpazo} or {times} if preferred
\usepackage{enumitem}
\usepackage{titlesec}
\usepackage[hidelinks]{hyperref}
\usepackage{xcolor}
\hypersetup{colorlinks=true, urlcolor=NavyBlue, linkcolor=NavyBlue}
\definecolor{NavyBlue}{RGB}{0,80,160}
\usepackage{needspace}
\usepackage{setspace}
\usepackage{etoolbox}
\pagestyle{empty}
\setlength{\parindent}{0pt}

% ================= PAGE MANAGEMENT =================
% (a) Fill each page: \flushbottom lets TeX stretch/shrink the elastic
%     ("plus/minus") spacing defined below so every page bottom aligns.
\raggedbottom
% (b) Never leave a section heading alone at the bottom of a page:
%     require at least 4 lines of room before starting a section.
\pretocmd{\section}{\Needspace*{4\baselineskip}}{}{}
% (c) Density knob: one number controls how tight the CV is.
%     0.8 = compact, 1.0 = normal, 1.3 = airy. Recompile after changing.
\newcommand{\density}{1.0}
% (d) Global line spacing: the biggest lever for filling the last page.
%     1.00 = tight, ~1.10 = airy. Tune together with the font size (10/11pt).
\setstretch{1.02}
\newlength{\gapS}\setlength{\gapS}{3pt plus 2pt minus 1pt}  % between list items
\newlength{\gapM}\setlength{\gapM}{4pt plus 2pt minus 1pt}  % between entries
\setlength{\gapS}{\density\gapS}\setlength{\gapM}{\density\gapM}
% ====================================================

% ---------- section heading style ----------
\titleformat{\section}{\large\bfseries\scshape}{}{0pt}{}[\vspace{-0.6em}\rule{\textwidth}{0.5pt}]
\titlespacing*{\section}{0pt}{1.1em plus 1.2em minus 0.4em}{0.6em plus 0.3em minus 0.2em}

% ---------- helper macros ----------
% \entry{Institution}{Date}{Role (italic)}{Place (italic)}
\newcommand{\entry}[4]{%
  \begin{samepage}%
  \textbf{#1} \hfill #2\par\nopagebreak
  {\itshape #3} \hfill {\itshape #4}\par\nopagebreak
  \end{samepage}}

% \pubitem{Authors}{Title}{Venue}{Links}   (leave a field empty if unused)
\newcommand{\pubitem}[4]{%
  \item #1. \textit{#2}. #3.\ifx\relax#4\relax\else\ #4\fi}

% \me  -> your bolded name in author lists
\newcommand{\me}{\textbf{Wenhan Cao}}
\newcommand{\mestar}{\textbf{Wenhan Cao}$^*$}

% [Paper]-style link
\newcommand{\plink}[2]{[\href{#2}{#1}]}

% award line:  \award{Text}{Year}
\newcommand{\award}[2]{#1 \hfill #2\par\vspace{2pt plus 2pt}}

% publication list environment (hanging indent, tight spacing)
% numbered list for publications: [1], [2], ...
\newenvironment{publist}[1][{[\arabic*]}]{%
  \begin{enumerate}[leftmargin=2.4em,labelsep=0.6em,label={#1},align=left,
                    itemsep=\gapS,parsep=0pt,topsep=3pt plus 1pt]%
  \interlinepenalty=10000\relax}% no page break inside a single reference
  {\end{enumerate}}

% ======================================================================
\newcommand{\acceptedpending}{accepted; volume and pages pending}
\newcommand{\earlyaccess}{Early Access}
\newcommand{\iclrstatus}{Under review at ICLR 2027}
\newcommand{\tmecstatus}{Under review at IEEE/ASME Transactions on Mechatronics (\textbf{T-MEC}), 2026}
\newcommand{\automaticastatus}{Under review at Automatica}
\begin{document}

% ---------- header ----------
\begin{center}
  {\Huge\bfseries Wenhan Cao}\\[6pt]
  \href{https://www.wenhancao.com}{www.wenhancao.com} \;|\;
  \href{mailto:wenhan_cao@outlook.com}{wenhan\_cao@outlook.com} \;|\;
  \href{https://scholar.google.com/citations?user=43xAy7MAAAAJ&hl=en}{Google Scholar}
\end{center}

% ---------- research interests ----------
\section{Research Interests}
My research develops a \textbf{probabilistic foundation for physical AI}: embodied systems, such as robots and autonomous vehicles, that must perceive, decide, and act under uncertainty. I design learning, estimation, and control methods that turn uncertainty from a nuisance into a first-class computational object. My long-term goal is to make probability the common language of reliable autonomy---enabling physical AI to reason about what it knows, adapt to what it doesn't, and earn human trust.

% ---------- work experience ----------
\section{Work Experience}
\entry{National University of Singapore}{August 2025 -- Present}
      {Eric and Wendy Schmidt AI in Science Postdoctoral Fellow,\\ Department of Electrical and Computer Engineering}{Singapore}
Supervisor: \href{https://sites.google.com/view/lzhao}{Prof.\ Lin Zhao}

% ---------- education ----------
\section{Education}
\entry{Tsinghua University}{September 2019 -- June 2025}
      {Ph.D., School of Vehicle and Mobility}{Beijing, China}
Supervisors: \href{https://scholar.google.com/citations?user=Dxiw1K8AAAAJ&hl=en}{Prof.\ Shengbo Eben Li} \&
\href{http://www2.coe.pku.edu.cn/faculty/liuchang/}{Prof.\ Chang Liu}
\vspace{\gapM}

\entry{University of Manchester}{January 2023 -- July 2024}
      {Visiting Ph.D.\ Researcher in Computer Science}{Manchester, UK}
Supervisor: \href{https://panweihit.github.io/}{Prof.\ Wei Pan}
\vspace{\gapM}

\entry{Technical University of Munich}{September 2023 -- December 2023}
      {Research Visit in Optimization and Control}{Munich, Germany}
Supervisor: \href{https://www.ce.cit.tum.de/en/itr/hirche/}{Prof.\ Sandra Hirche}
\vspace{\gapM}

\entry{Beijing Jiaotong University}{September 2015 -- June 2019}
      {B.Eng.\ in Electrical Engineering}{Beijing, China}
GPA ranking: 1/305

% ---------- publications ----------
\section{Selected Publications}
{\small\itshape $^*$ denotes equal contribution. A full list is available in the
\href{https://www.wenhancao.com/data/Publication_list_by_type.pdf}{publication list} or on
\href{https://scholar.google.com/citations?hl=en&user=43xAy7MAAAAJ}{Google Scholar}.}

\begin{publist}

\item \me, Xuyang Chen, Shuyuan Wang \& Lin Zhao. \textit{How Much Information is Needed for Accurate Kalman Filtering?} In Advances in Neural Information Processing Systems (\textbf{NeurIPS}), 2026. \plink{Code}{https://anonymous.4open.science/r/rdkf-4B31}

\pubitem{Shiqi Liu$^*$, \mestar, Chang Liu, Zeyu He, Tianyi Zhang, Yinuo Wang \& Shengbo Eben Li}
        {One Filters All: A Generalist Filter For State Estimation}
        {In Advances in Neural Information Processing Systems (\textbf{NeurIPS}), 38 (Main Conference):21276--21302, 2025}
        {\plink{Paper}{https://openreview.net/forum?id=EGK487IYAW}
         \plink{Poster}{https://neurips.cc/virtual/2025/poster/119149}}

\pubitem{\me, Alexandre Capone, Rishabh Dev Yadav, Sandra Hirche \& Wei Pan}
        {Computation-Aware Learning for Stable Control with Gaussian Process}
        {In Robotics: Science and Systems (\textbf{RSS}), 2024}
        {\plink{Paper}{https://www.roboticsproceedings.org/rss20/p004.html}
         \plink{Poster}{https://drive.google.com/file/d/1rSKLy-SEqgMR7Z8Ius9dhHuqOROflWe0/view?usp=sharing}
         \plink{Recording}{https://www.youtube.com/watch?v=AE0iPGlVM30}}

\pubitem{\me\ \& Wei Pan}
        {Impact of Computation in Integral Reinforcement Learning for Continuous-Time Control}
        {In International Conference on Learning Representations (\textbf{ICLR}), 2024. \textbf{(Spotlight)}}
        {\plink{Paper}{https://openreview.net/forum?id=xJEd8PkdNz}
         \plink{Poster}{https://drive.google.com/file/d/16z45LOEU8xaPaEoFFIc2hfIneyXCCio_/view?usp=sharing}
         \plink{Code}{https://github.com/anonymity678/Computation-Impacts-Control}}

\pubitem{\me, Tianyi Zhang, Zeju Sun, Chang Liu, Stephen S.-T. Yau \& Shengbo Eben Li}
        {Nonlinear Bayesian Filtering with Natural Gradient Gaussian Approximation}
        {IEEE Transactions on Pattern Analysis and Machine Intelligence (\textbf{TPAMI}), 48(7):8581--8597, 2026}
        {\plink{Paper}{https://arxiv.org/abs/2410.15832}
         \plink{Code}{https://github.com/wenhancao/NANO-idlab}
         \plink{Slides}{https://drive.google.com/file/d/1O8qNHLxzhbikJ_4Ae-LQPQEzaFJre81U/view?usp=sharing}}

\pubitem{\me, Shiqi Liu, Chang Liu, Zeyu He, Stephen S.-T. Yau \& Shengbo Eben Li}
        {Convolutional Bayesian Filtering}
        {IEEE Transactions on Automatic Control (\textbf{TAC}), pp. 1--8, 2026}
        {\plink{Paper}{https://arxiv.org/abs/2404.00481}
         \plink{Slides}{https://drive.google.com/file/d/1vuGl7mLdf4Muuaur2buQaIbgiV8pHmTu/view}}

\pubitem{\me, Chang Liu, Zhiqian Lan, Shengbo Eben Li, Wei Pan \& Angelo Alessandri}
        {Robust Bayesian Inference for Moving Horizon Estimation}
        {\textbf{Automatica}, 173, 112108, 2025}
        {\plink{Paper}{https://www.sciencedirect.com/science/article/abs/pii/S0005109824006010}
         \plink{Code}{https://github.com/wenhancao/robust-Bayesian-inference-for-MHE.git}}

\pubitem{Jingliang Duan, \me, Yang Zheng \& Lin Zhao}
        {On the Optimization Landscape of Dynamic Output Feedback Linear Quadratic Control}
        {IEEE Transactions on Automatic Control (\textbf{TAC}), 69(2):920--935, 2024}
        {\plink{Paper}{https://ieeexplore.ieee.org/abstract/document/10124022}
         \plink{Code}{https://github.com/soc-ucsd/LQG_gradient/tree/master/dLQR}}

\pubitem{Kangjie Zhou, Zhejia Wen, Zhiyong Zhuo, Zike Yan, Pengying Wu, Shuaiyang Li, Han Gao, Kang Ding, \me, Wei Pan \& Chang Liu}
        {CoINS: Counterfactual Interactive Navigation via Skill-Aware VLM}
        {IEEE Transactions on Automation Science and Engineering (\textbf{T-ASE}), 2026}
        {\plink{Paper}{https://arxiv.org/pdf/2601.03956}
         \plink{Project}{https://coins-internav.github.io/}}

\pubitem{Tianyi Zhang, \me, Chang Liu, Tao Zhang, Jiangtao Li \& Shengbo Eben Li}
        {Robust State Estimation for Legged Robots with Dual Beta Kalman Filter}
        {IEEE Robotics and Automation Letters (\textbf{RA-L}), 10(8):7955--7962, 2025}
        {\plink{Paper}{https://arxiv.org/abs/2411.11483}}

\pubitem{Shiqi Liu, \me, Chang Liu, Tianyi Zhang \& Shengbo Eben Li}
        {Convolutional Unscented Kalman Filter for Multi-Object Tracking with Outliers}
        {IEEE Transactions on Intelligent Vehicles (\textbf{T-IV}), 2024}
        {\plink{Paper}{https://ieeexplore.ieee.org/document/10643343}}

\pubitem{Tianyi Zhang, \me, Chang Liu, Feihong Zhang, Wei Wu \& Shengbo Eben Li}
        {NANO-SLAM: Natural Gradient Gaussian Approximation for Vehicle SLAM}
        {In IEEE International Conference on Intelligent Transportation Systems (\textbf{ITSC}), pp. 4014--4019, 2025}
        {\plink{Paper}{https://arxiv.org/abs/2504.19195}}

\pubitem{Tianyi Zhang, \me\ \& Shengbo Eben Li}
        {Natural Gradient Gaussian Approximation Filter with Positive Definiteness Guarantee}
        {In American Control Conference (\textbf{ACC}), pp. 3224--3229, 2026}
        {\relax}

\pubitem{\me, Chang Liu, Zhiqian Lan, Yingxi Piao \& Shengbo Eben Li}
        {Generalized Moving Horizon Estimation for Nonlinear Systems with Robustness to Measurement Outliers}
        {In American Control Conference (\textbf{ACC}), 2023}
        {\plink{Paper}{https://ieeexplore.ieee.org/document/10156391}
         \plink{Code}{https://github.com/wenhancao/robust-Bayesian-inference-for-MHE.git}
         \plink{Slides}{https://drive.google.com/file/d/1I4Ihmdwlf3kf3sYs4YXGwT6WWOHkr-Pd/view?usp=sharing}}

\pubitem{\me, Jingliang Duan, Shengbo Eben Li, Chen Chen, Chang Liu \& Yu Wang}
        {Primal-Dual Estimator Learning Method with Feasibility and Near-Optimality Guarantees}
        {In IEEE Conference on Decision and Control (\textbf{CDC}), 2022}
        {\plink{Paper}{https://ieeexplore.ieee.org/document/9992814}
         \plink{Slides}{https://drive.google.com/file/d/1Nueg0BkHXsr6s7BwFkDdz3MGJXMRvt9i/view?usp=sharing}}

\pubitem{\me, Jianyu Chen, Jingliang Duan, Shengbo Eben Li, Yao Lyu, Ziqing Gu \& Yuhang Zhang}
        {Reinforced Optimal Estimator}
        {In Modeling, Estimation and Control Conference (\textbf{MECC}), 2021. \textbf{(Student Best Paper Finalist)}}
        {\plink{Paper}{https://www.sciencedirect.com/science/article/pii/S240589632102245X}
         \plink{Slides}{https://drive.google.com/file/d/1ZXlsONIwnuPpb4IFkeCRYL7clDx7hTor/view?usp=sharing}}

\end{publist}


% ---------- preprints ----------
\section{Preprints}
\begin{publist}[{[P\arabic*]}]

\pubitem{\me, Keyu Yan, Yang Feng, Rick Siow Mong Goh \& Jun Zhou}
        {HIRE: Hippocampal Replay for Test-Time Agentic Reinforcement Learning}
        {\iclrstatus}
        {}

\pubitem{Keyu Yan$^*$, \mestar, Shen'ao Wang, Yujie Yang, Wu Junke, Xuyang Chen \& Lin Zhao}
        {Generating Parallel Worlds of a Flight}
        {\iclrstatus}
        {}

\pubitem{Keyu Yan, Xuyang Chen, \me\ \& Lin Zhao}
        {Evaluating OOD Actions in Offline RL via Quantile Advantage}
        {\iclrstatus}
        {\plink{Code}{https://anonymous.4open.science/r/adac-14D0}}

\pubitem{\me, Keyu Yan, Kuangyi Wang, Xuyang Chen \& Lin Zhao}
        {Trajectory Editing as Posterior Inference for Offline Reinforcement Learning}
        {\iclrstatus}
        {\plink{Code}{https://anonymous.4open.science/r/bite-055D}}

\pubitem{Xuyang Chen, Keyu Yan, Danze Chen, Hongzong Li, \me\ \& Lin Zhao}
        {Learning Robotic Skills from Offline Preference-based RL}
        {\iclrstatus}
        {\plink{Code}{https://anonymous.4open.science/r/PSAC-6686}}

\pubitem{Rui Huang, Xin Chen, \me, Lidong Li, Yizhe Xiao, Yichao Gao, Yuqi Shi \& Lin Zhao}
        {Agile Modular Aerial Robot Load Transport with Arbitrary Configurations}
        {\tmecstatus}
        {\relax}

\pubitem{Han Gao, \me, Ieng Hou U, Kangjie Zhou, Ying Sun, Angelo Alessandri \& Chang Liu}
        {Chance-Constrained Gaussian Filtering}
        {\automaticastatus}
        {\relax}

\end{publist}


% ---------- honors ----------
\section{Honors \& Awards}
\award{Eric and Wendy Schmidt AI in Science Postdoctoral Fellowship}{2026}
\award{Study Abroad Fund, Tsinghua University}{2022}
\award{Student Best Paper Finalist, Modeling, Estimation and Control Conference}{2021}
\award{China National Scholarship}{2016}
\award{First Prize Scholarship, Beijing Jiaotong University}{2016, 2017 \& 2018}

% ---------- software ----------
\section{Software}
\textbf{General Optimal control Problems Solver (GOPS)} --- an easy-to-use reinforcement learning (RL) solver package designed to build real-time, high-performance controllers for industrial applications. I was primarily responsible for the core design and implementation of the \textbf{trainer}, \textbf{sampler}, and \textbf{buffer} modules.
\plink{Docs}{https://gops.readthedocs.io/en/latest/index.html}
\plink{Code}{https://github.com/Intelligent-Driving-Laboratory/GOPS}
\plink{Paper}{https://www.sciencedirect.com/science/article/pii/S2772424723000070}

% ---------- projects ----------
\section{Selected Research Projects}
\textbf{RL Post-Training of MindGPT} \hfill November 2024\par
{\itshape Project Member, Internship at \href{https://www.lixiang.com/}{Li Auto}}
\vspace{\gapM}

\textbf{Reinforcement Learning for Autonomous Driving Decision and Control on Open Roads} \hfill May 2024\par
{\itshape Project Member} --- Supported by \href{https://global.toyota/en/}{Toyota} \& \href{https://www.idriverplus.com/}{IdriverPlus}.
\vspace{\gapM}

\textbf{State Estimation for Warehouse Automated Logistics Vehicles} \hfill May 2022\par
{\itshape Student Project Leader} --- Supported by \href{https://www.geekplus.com/en/}{Geek+}; the technology has been implemented in the company's product lineup.
\vspace{\gapM}

\textbf{Networked Modeling and Cooperative Control of Connected and Automated Electric Vehicles}\par
{\itshape Student Project Leader} --- Supported by MOST. \hfill November 2020

% ---------- talks ----------
\section{Invited Talks}
\begin{itemize}[leftmargin=1.4em,itemindent=-1.4em,labelsep=0pt,label={},itemsep=\gapS,parsep=0pt,topsep=3pt plus 1pt]\interlinepenalty=10000\relax
\item \textit{NANO Filter: Bayesian Filtering with Natural Gradient Gaussian Approximation}. Department of Astronomy, Tsinghua University, Beijing, China, hosted by Prof.\ Zheng Cai, August 2024.
\item \textit{Convolutional Bayesian Filtering}. Department of Mathematical Sciences, Tsinghua University, Beijing, China, hosted by Prof.\ Stephen S.-T. Yau, February 2024.
\item \textit{Learning-based State Estimation Methods}. Technical University of Munich, Munich, Germany (online), hosted by Prof.\ Sandra Hirche, February 2023.
\end{itemize}

% ---------- teaching ----------
\section{Teaching Experience}
\textbf{Tsinghua University, School of Vehicle and Mobility} \hfill February 2022 -- May 2022\par
{\itshape Teaching Assistant, 150051 Intelligent Vehicles}

% ---------- services ----------
\section{Professional Services}
\textbf{Conference Reviewer:} NeurIPS, ICML, ICLR, AAMAS, L4DC, CDC, ACC \& IFAC NMPC\par
\vspace{2pt}
\textbf{Journal Reviewer:} TPAMI, T-RO, T-RL, TAC, Automatica, RA-L, T-ITS, T-ASE, TII \& TNNLS

\end{document}
